\documentclass[11pt]{article}
\usepackage[T1]{fontenc}
\usepackage[utf8]{inputenc}
\usepackage{amsmath,amssymb,amsthm}
\usepackage[margin=1in]{geometry}
\usepackage{booktabs}
\usepackage[hidelinks,pdftitle={Baryogenesis in the Psion Hypothesis},pdfauthor={Karlos Marden Maia}]{hyperref}

\newtheorem{theorem}{Theorem}
\newtheorem{proposition}[theorem]{Proposition}
\newtheorem{corollary}[theorem]{Corollary}
\newtheorem{postulate}{Postulate}
\theoremstyle{remark}
\newtheorem{remark}{Remark}

\newcommand{\Lm}{\Lambda_{\rm meta}}
\newcommand{\Np}{N_\psi}

\title{Baryogenesis in the Psion Hypothesis:\\
Psion-Number Conservation, No-Go Results below the Condensation Scale,\\
and Level Bifurcation as a Candidate Mechanism}
\author{Karlos Marden Maia\\ \small Independent researcher, Marietta, Georgia, USA}
\date{October 2026}

\begin{document}
\maketitle

\begin{abstract}
The Psion Hypothesis (PH) builds all matter from a single six-label fermion, the psion, interacting through one four-fermion Lagrangian. We examine what this structure implies for the baryon asymmetry of the Universe. Four results are exact. (i) Psion number is conserved by every term of the Lagrangian, and the continuous internal symmetry is only $U(1)_N\times U(1)_M$. (ii) Field-theory CP is exact. (iii) Observed reactions select a unique ontology for antimatter, composites of antipsions, and require exactly two additive charges, $B$ and $L$, with $\Np/3=B+L$; the PH therefore conserves $B+L$ and allows only $B-L$ to change, the reverse of the Standard Model. (iv) The Lagrangian supplies $B$ neither as a Noether charge nor as a topological charge, and the cascade level is not an internal degree of freedom. On this basis we show, with pre-registered numerical tests, that every mechanism operating below the condensation scale $\Lm\simeq1.2\times10^{12}$~GeV fails. We then formulate a candidate mechanism at the condensation scale, \emph{level bifurcation}: each psion triple binds into a baryonic or a leptonic rung with a CP-odd bias, and the rung is conserved afterwards. If the mechanism holds, $B=-L$, charge neutrality and a final state of hydrogen plus two net antineutrinos per baryon follow from its symmetry alone, and the efficiency reaches $\eta_B/\varepsilon\simeq3.6\times10^{-3}$ without fine tuning. The mechanism rests on three conditions that are stated, not derived: a conserved rung label, a CP bias of global sign, and the departure from equilibrium and dilution of the condensation transition. The PH thus makes structural predictions independent of the mechanism: no nucleon decay with $\Delta B=\Delta L$, no $0\nu\beta\beta$, no $n$--$\bar n$ oscillation, and, if the asymmetry formed with zero net psion number, a cosmic antineutrino excess of two per baryon.
\end{abstract}

\section{Introduction}

A dynamical origin of the baryon asymmetry, $\eta_B=n_B/n_\gamma=6.1\times10^{-10}$ \cite{planck}, requires baryon-number violation, C and CP violation, and departure from thermal equilibrium \cite{sakharov}. Successful frameworks, such as electroweak baryogenesis \cite{krs}, leptogenesis \cite{fy,bdp} and the Affleck--Dine mechanism \cite{ad}, derive all three ingredients within a specified model.

The Psion Hypothesis \cite{phf} is a single-Lagrangian framework in which the observed fermions are three-psion composites on a mass cascade $m(\ell)=\Lm e^{-\ell/b_e}$. Earlier versions assigned the asymmetry to a chirality parameter $\varepsilon_\psi=25/169$ combined with a Standard-Model sphaleron conversion. Here we take the Lagrangian at face value and ask what it permits. The answer is restrictive. The constraints turn out to be strong enough to fix the \emph{form} a viable mechanism must take, and we formulate and test one such mechanism.

The paper proceeds from exact statements to conditional ones. Section~\ref{sec:fw} recalls the framework. Section~\ref{sec:exact} proves the conservation structure. Section~\ref{sec:nogo} closes the mechanisms available below $\Lm$. Section~\ref{sec:bif} formulates level bifurcation and states the conditions under which it accounts for $\eta_B$. Section~\ref{sec:pred} collects the predictions. Every numerical test was pre-registered: the definition, with a SHA-256 hash and timestamp, was fixed before execution, and all outputs, including failed controls, are public \cite{repo}.

\section{Framework}\label{sec:fw}

The Lagrangian is
\begin{equation}\label{eq:L}
\mathcal L=\bar\psi\, i\gamma^\mu\partial_\mu\psi-G_{\rm UV}(\bar\psi\psi)^2-(\bar\psi J^a\psi)(\bar\psi\Sigma^a\psi),
\end{equation}
where $\psi$ carries a six-valued label $k=0,\dots,5$ (internal spin $j=5/2$, $m_k=k-5/2$). $J^a$ and $\Sigma^a$ are the SU(2)$_{\rm meta}$ generators restricted to the left ($k\le2$) and right ($k\ge3$) sectors. The labels carry charges $Q_k=(1,-5,-5,10,4,4)/9$. Composites of three distinct labels (``triples'') number 20. Among them are the electron $(0,1,2)$, the proton $(0,4,5)$ and the neutron $(0,1,5)$. The cascade level $\ell$ is an independent quantum number: the nucleon scale is $\ell=8$ (610~MeV) and the electron $\ell=10$ \cite{phf}. Three neutral triples are listed in \cite{phf} as dual-sector, each being either a baryon or a neutrino: $(0,1,5)=n/\nu_\tau$, $(0,2,4)=\Sigma^0/\nu_\mu$, $(1,2,3)=\Lambda^0/\nu_e$.

\section{Exact conservation structure}\label{sec:exact}

\begin{theorem}[Psion number]\label{th:N}
$\Np$, the charge of $\psi\to e^{i\alpha}\psi$, is conserved by every term of \eqref{eq:L}. In the condensed phase it is the only exact additive charge.
\end{theorem}
\noindent The first statement is immediate. The second follows from an exact charge audit over the 20 triples and four condensate variants, which leaves no further candidate \cite{repo}.

\begin{theorem}[Continuous internal symmetry]\label{th:sym}
The Lie algebra of transformations $\psi\to e^{tX}\psi$ that leave the interaction in \eqref{eq:L} invariant is two-dimensional and abelian. It is generated by the identity and by ${\rm diag}(m_k)$, i.e.\ the group is $U(1)_N\times U(1)_M$ with $M=\sum m_k$. No continuous chiral extension ($X\otimes\gamma_5$) survives the $G_{\rm UV}$ term. The winding condensate $J^+\Sigma^++J^-\Sigma^-$ breaks $U(1)_M$ and preserves $U(1)_N$.
\end{theorem}
\noindent The invariance condition is linear in $X$. Its solution space was computed exactly in the 36-dimensional space of anti-Hermitian $6\times6$ matrices.

\begin{proposition}[CP]\label{pr:cp}
The coupling tensors of \eqref{eq:L} are real and Hermitian, so field-theory CP is exact. The chirality parameter $\varepsilon_\psi$ is even in any hopping phase, hence CP-even. The six Z6 vacua violate CP pairwise, but conjugate vacua $(k,6-k)$ are degenerate.
\end{proposition}

\begin{theorem}[Antimatter and the observed charges]\label{th:ont}
Consider the two readings of antimatter: (O$^+$) a triple of psions with inverted labels $k\to(6-k)\bmod6$; (O$^-$) the same triple of antipsions, with $\Np=-3$ and charge $-Q$. Test both against the observed reactions $n\to pe^-\bar\nu_e$, $p\to ne^+\nu_e$, $e^-p\to n\nu_e$, $e^+e^-\to2\gamma,3\gamma$, $\gamma\to e^+e^-$, $p\bar p\to\pi^+\pi^-\pi^0$, $\pi^0\to2\gamma$, $\pi^+\to\mu^+\nu_\mu$ and $\mu^-\to e^-\bar\nu_e\nu_\mu$. Reading O$^+$ conserves neither $\Np$ nor charge and assigns $\bar\Sigma^+$ the charge $+1$. Reading O$^-$ passes all tests, with $\Np=0$ for every boson.
\end{theorem}

\begin{theorem}[Charges required by observation]\label{th:BL}
Treat the cascade class as a label: $H$ (level 8: $p,n,\Lambda,\Sigma^0$), $E$ (level 10 and its excitations: $e,\mu$), $\nu$ (three neutrinos). The additive charges conserved by the observed reactions $\beta^\pm$ decay, $e^-$ capture, $\pi^+\to\mu^+\nu_\mu$, $\mu$ decay, $\nu_\mu n\to\mu^-p$, $\bar\nu_ep\to e^+n$, $\Lambda\to p\pi^-,n\pi^0$, $\Sigma^0\to\Lambda\gamma$ and neutrino oscillations span exactly a two-dimensional space. Its basis is $B$ (number of $H$ states) and $L$ (number of $E$ and $\nu$ states), and $\Np/3=B+L$.
\end{theorem}

\begin{corollary}
Since $\Np$ is exact, $B+L$ is conserved. Baryogenesis in the PH requires a change in $B-L$, the reverse of the Standard Model, where $B-L$ is conserved and sphalerons change $B+L$. The Standard-Model sphaleron conversion therefore does not apply.
\end{corollary}

\begin{theorem}[Baryon number is not available in \eqref{eq:L}]\label{th:noB}
(a) $B$ is not in the span of the Noether charges $(M,Q,\Np)$ evaluated on the observed states. (b) The vacuum manifold of the condensed phase has $\pi_2=\pi_3=0$, so no point-like topological charge exists and there is no analogue of the Skyrme baryon number \cite{witten}. (c) Composite states are labelled by their triple alone; the cascade level carries no internal quantum number.
\end{theorem}

\begin{corollary}[Dual-sector instability]\label{co:dual}
With only $\Np$ and $Q$ conserved, the processes $p\to\nu\pi^+$, $n\to\nu\pi^0$, $n\to3\nu$, $pe^-\to\nu_\mu\nu_\tau$ and the same-triple level jumps $n\to\nu_\tau\gamma$, $\Lambda^0\to\nu_e\gamma$, $\Sigma^0\to\nu_\mu\gamma$ are all allowed, and all have $\Delta B=-\Delta L$. The measured limits are $p\to\bar\nu\pi^+>3.9\times10^{32}$~yr and $n\to\bar\nu\pi^0>1.1\times10^{33}$~yr \cite{sk}, and $n\to{\rm invisible}>5.8\times10^{29}$~yr \cite{kamland}. They require an amplitude suppression $\lesssim10^{-8}$ relative to a contact operator at $\Lm$. Same-triple jumps, which no label-changing selection rule of \eqref{eq:L} can forbid, require $\lesssim5\times10^{-30}$ relative to an electromagnetic transition (reference lifetime $10^{29}$~yr). A conserved baryon label must therefore be added to the PH.
\end{corollary}

\noindent The count of neutral triples is five, not three: $(0,1,4)$ and $(0,2,5)$ are also neutral. They carry the residual discrete charge $D=(2M-6Q+3)\bmod4=2$, which is exact in the condensed phase.

\section{No-go results below the condensation scale}\label{sec:nogo}

\paragraph{Emergent sphalerons.} Any sphaleron-like configuration of the PH is built from psions and conserves $\Np$, hence $B+L$. It cannot convert a lepton asymmetry into baryons in the Standard-Model manner.

\paragraph{Temperature histories.} CPT gives particle and antiparticle equal masses and equal equilibrium abundances. With exact CP and a symmetric initial state, a Boltzmann network of the cascade levels returns zero asymmetry for any programme of cooling and reheating; reheating only dilutes an existing asymmetry. Temperature controls the composition across levels, not the matter--antimatter ratio.

\paragraph{Z6 domains with spontaneous CP.} Inside a CP-violating domain, assembly of composites out of equilibrium does generate an asymmetry, $\Delta p\propto\varepsilon$, together with an equal antipsion reservoir. However, the Z6 walls form at $T\approx3$~MeV, when the assembly rate is $\Gamma/H\sim10^{-37}$. A domain's maximal comoving size today is $\sim30$~pc, against the $\gtrsim1$~Gpc required of antimatter regions \cite{cdg}. And degenerate conjugate vacua give a zero global average.

\paragraph{CP terms built from PH operators.} We tested 16 CP-odd patterns built from the PH's own operators (commutators of $Q$, $M$, $n_R$, the Z6 index and the energy with $h$ and $J\!\cdot\!\Sigma$, and products in rearrangements), singly and in pairs. Fitted to the observed composition, none performs better than random directions. Inserting a phase directly, $J\!\cdot\!\Sigma\to{\rm Re}(e^{-i\varphi}O)$, with full density-matrix assembly kinetics, generates an asymmetry with $\varphi\sim10^{-4}$--$10^{-2}$. But exotic triples receive 2--2500 times the proton asymmetry, and the electron asymmetry has the wrong sign in most runs. Without a conserved label, no CP source selects the observed composites.

\paragraph{Minimal extension.} A two-valued species index $\psi\to\psi_s$, $s\in\{B,L\}$, with currents summed over species and $SU(2)_{\rm species}$ broken, makes $B$ and $L$ Noether charges. It allows all observed reactions and forbids the processes of Corollary~\ref{co:dual} exactly. It preserves the PH's numerical outputs only if the index does not enter the one-loop coefficient $|\beta_0|=N_{\rm Dirac}+2/N_f$: doubling $N_f$ gives $m_e=0.065$~MeV instead of 0.509~MeV. If $B-L$ is to be violated by species mixing, the proton lifetime bounds the mixing coefficient to $c\le1.3\times10^{-8}$. Mixing is then active only above $2.7\times10^{15}$~GeV, outside the range of the contact description.

\begin{proposition}\label{pr:window}
Within \eqref{eq:L} and its minimal extension, no mechanism operating below $\Lm$ produces the baryon asymmetry. The asymmetry must be fixed at the condensation scale or above, or be an initial condition.
\end{proposition}

\section{Level bifurcation at formation}\label{sec:bif}

\subsection{Postulates}

The dual sector, which causes the instability of Corollary~\ref{co:dual}, also supplies the structure that the failed mechanisms lacked: two final channels from a common initial state. That is the minimal setting in which CP violation is compatible with CPT and unitarity \cite{kw}. We formulate the mechanism as three postulates.

\begin{postulate}[Rung label]\label{P1}
A psion triple formed at the condensation scale binds into the baryonic rung ($t_B$) or the leptonic rung ($t_L$). After formation, the rung is conserved by all processes.
\end{postulate}

\begin{postulate}[Biased bifurcation]\label{P2}
The two channels are populated with a CP-odd bias $\varepsilon$, of opposite sign for antipsion triples and of the same sign throughout the observable Universe. The CP source is $\varepsilon$ times the net formation rate, so it vanishes in equilibrium; dissociation followed by re-formation washes out the channel imbalance.
\end{postulate}

\begin{postulate}[Departure from equilibrium]\label{P3}
Formation occurs during the condensation transition with a finite departure from equilibrium, parametrised by the contact-rate normalisation $c_\Gamma$. It is followed by an entropy dilution $D\le1$.
\end{postulate}

\noindent Postulate~\ref{P1} is the conserved label required independently by Corollary~\ref{co:dual}. The mechanism therefore introduces no ingredient beyond what nucleon stability already demands, apart from the bias.

\subsection{Kinetics}

For each triple $t$ and formation rate $k_t$, the matter equations read
\begin{align}
\frac{dt_B}{dx}&\ni \tfrac12 k_tP_t-\frac{k_t}{K}t_B+\tfrac12\varepsilon s_tJ_t, &
\frac{dt_L}{dx}&\ni \tfrac12 k_tP_t-\frac{k_t}{K}t_L-\tfrac12\varepsilon s_tJ_t,
\end{align}
with $P_t$ the product of the three psion abundances and $J_t=k_t[P_t-(t_B+t_L)/K]$ the net formation rate. The antimatter equations are the same with $\varepsilon\to-\varepsilon$. The network also contains transitions within each rung, rearrangements $ab\to cd$ within each rung, annihilation $t_r\bar t_r\to$ radiation, and $\psi\bar\psi$ annihilation; the rates and binding weights are those of the assembly network of the PH audit \cite{repo}. We take $s_t=1$, a bias of uniform orientation.

\begin{proposition}[Structural consequences]\label{pr:struct}
For a symmetric initial state, Postulates~\ref{P1}--\ref{P2} imply, at all times,
\begin{equation}
B=-L,\qquad Q_{\rm tot}=0,\qquad \Delta\Np=0 .
\end{equation}
If $p$, $e^-$ and the neutrinos are the lightest states of their sectors, the late-time state is hydrogen plus two net antineutrinos per baryon, whatever the composition at formation.
\end{proposition}
\begin{proof}
The map $(\psi,t_B,t_L,\bar t_B,\bar t_L)\to(\bar\psi,\bar t_L,\bar t_B,t_L,t_B)$ leaves the equations invariant: it exchanges $+\varepsilon$ with $-\varepsilon$ and pairs $t_B\bar t_B$ with $\bar t_Lt_L$ in the annihilation terms. It also leaves the symmetric initial state invariant. Hence $t_B=\bar t_L$, $t_L=\bar t_B$ and $\psi=\bar\psi$ at all times, which gives $B=\sum(t_B-\bar t_B)=-\sum(t_L-\bar t_L)=-L$ and $\Delta\Np=0$. The charge of each triple enters $B$ and $L$ with opposite signs, so $Q_{\rm tot}=0$. With $B$, $L$ and $Q$ conserved, baryonic excesses decay to protons and leptonic ones to $e^-$ and neutrinos. Neutrality gives $N_e=N_p=B$, and $L=N_e-(N_{\bar\nu}-N_\nu)=-B$ gives $N_{\bar\nu}-N_\nu=2B$.
\end{proof}

\noindent Proposition~\ref{pr:struct} removes the composition problem of Section~\ref{sec:nogo}. Exotic triples need not be absent at formation; they decay into the observed states because the labels are conserved. We verified the proposition numerically in a 92-species network. In all runs with $|B|$ above numerical noise, $|B+L|/|B|\le6\times10^{-10}$, $|Q_{\rm tot}|\lesssim10^{-14}$ and $|\Delta\Np|\lesssim10^{-13}$. A bias of $\varepsilon=0.1$ produces no asymmetry from an equilibrium initial state ($\lesssim10^{-16}$), and the response is linear in $\varepsilon$ to $3\times10^{-8}$.

\subsection{Efficiency}

Table~\ref{tab:eff} gives the efficiency $\eta_B/\varepsilon$ for $D=1$ as a function of $c_\Gamma$, for the two binding laws of the network.

\begin{table}[h]
\centering
\begin{tabular}{lcccccc}
\toprule
$c_\Gamma$ & $10^{-5}$ & $10^{-4}$ & $10^{-3}$ & $10^{-2}$ & $10^{-1}$ & 1 \\
\midrule
law L1 & $1.9\times10^{-4}$ & $1.4\times10^{-3}$ & $3.6\times10^{-3}$ & $3.6\times10^{-4}$ & $<10^{-12}$ & $<10^{-11}$ \\
law L2 & $1.9\times10^{-4}$ & $1.3\times10^{-3}$ & $2.6\times10^{-3}$ & $4.9\times10^{-4}$ & $<10^{-13}$ & $<10^{-11}$ \\
\bottomrule
\end{tabular}
\caption{Efficiency $\eta_B/\varepsilon$ of level bifurcation (uniform bias, $D=1$). For $c_\Gamma\ge0.1$ formation is in equilibrium and the asymmetry is washed out.}\label{tab:eff}
\end{table}

\noindent Between $c_\Gamma=10^{-5}$ and $10^{-2}$ the efficiency changes by a factor $\sim20$ and keeps its sign. The mechanism therefore does not depend on a finely tuned rate. A bias weighted by triple chirality, $s_t=(n_L-n_R)/3$, changes sign between the two binding laws; such a weighting cannot fix the sign of $\eta_B$.

\paragraph{Consistency condition.} The mechanism reproduces the observed value if
\begin{equation}\label{eq:cond}
\varepsilon\,D\simeq\frac{6.1\times10^{-10}}{\eta_B/\varepsilon}\in\left[1.7\times10^{-7},\,3.2\times10^{-6}\right]
\end{equation}
over the non-equilibrium range of Table~\ref{tab:eff}. This is a requirement on the physics of the transition, not a prediction. A natural scale for $|\varepsilon|$ is the chirality parameter $\varepsilon_\psi=25/169$; with that value, \eqref{eq:cond} requires a dilution $D\sim10^{-6}$--$10^{-5}$. Existing estimates of the condensation transition (Coleman--Weinberg with thermal inflation) give $D^{-1/3}=\xi^*\approx3.6$--$39$ \cite{repo}. Deriving $D$ is part of Condition~C3 below.

\subsection{Conditions}

The mechanism is identified and its structure verified, but its magnitude rests on modelled contact kinetics and three underived inputs.
\begin{itemize}
\item[(C1)] \emph{Rung label.} Derive a conserved baryon label (Postulate~\ref{P1}). The same principle must secure nucleon stability (Corollary~\ref{co:dual}).
\item[(C2)] \emph{Sign and size of the bias.} $\varepsilon_\psi$ is CP-even (Proposition~\ref{pr:cp}) and can at most set the size. The sign must come from a vacuum orientation common to the observable Universe; the natural candidate is the winding direction of the condensate, which carries a definite chirality through its Jackiw--Rebbi zero mode \cite{jr}.
\item[(C3)] \emph{Transition dynamics.} Derive the departure from equilibrium ($c_\Gamma$) and the dilution $D$ from the condensation transition.
\end{itemize}
If C1--C3 hold and satisfy \eqref{eq:cond}, the PH accounts for the baryon asymmetry and nucleon stability through a single principle.

\section{Predictions}\label{sec:pred}

The following predictions do not depend on the mechanism of Section~\ref{sec:bif}, unless stated otherwise.
\begin{enumerate}
\item No nucleon decay with $\Delta B=\Delta L$ ($p\to e^+\pi^0$, $p\to e^+\gamma$, \dots): these change $\Np$ and are forbidden exactly by Theorem~\ref{th:N}.
\item If any nucleon decay is observed, it has $\Delta B=-\Delta L$ ($p\to\nu\pi^+$, $n\to\nu\pi^0$, $n\to3\nu$).
\item No neutrinoless double-beta decay and no $n$--$\bar n$ oscillation: both change $\Np$.
\item If the asymmetry formed with zero net psion number (Section~\ref{sec:bif}): a cosmic antineutrino excess of two per baryon, i.e.\ a neutrino degeneracy parameter $|\xi_\nu|\approx2.5\times10^{-9}$ per flavour (shared over two flavours). This is far below current bounds, and it is the characteristic, if difficult, signature of the mechanism.
\end{enumerate}

\section{Discussion}

The PH inverts the Standard-Model structure of baryon and lepton number. This places it outside the families of mechanisms that rely on sphaleron reprocessing, and it rules out several routes that would otherwise be natural. The constraint that remains is sharp: the theory needs a conserved baryon label. Without it, the dual-sector identification makes nucleons unstable. With it, a CP-biased bifurcation at formation yields the observed composition with no tuning of the composition itself.

Two limitations bound the present results. First, the kinetic networks use contact rates and binding weights that are models, not derived amplitudes. The structural statements (Proposition~\ref{pr:struct}) are independent of these choices; the efficiencies of Table~\ref{tab:eff} are not. Second, Postulate~\ref{P1} is an addition to \eqref{eq:L}. Section~\ref{sec:nogo} shows that the simplest realisation, a non-running species label, is equivalent to postulating baryon-number conservation. Deriving the label is the central open problem.

A related tension concerns dark matter. If the top of the cascade were a thermal relic, $\Omega_{\rm DM}h^2=0.12$ would require a mass of $\sim56$~TeV (cascade level $\approx4.8$). The canonical PH value of $7.2\times10^{10}$~GeV would overclose the Universe; a non-thermal production route is therefore required.

\section{Conclusions}

Equation~\eqref{eq:L} conserves psion number and CP exactly. Observation then requires separate baryon and lepton numbers with $B+L$ fixed, and the PH must change $B-L$ to produce an asymmetry. Below the condensation scale this is impossible within the theory and its minimal extension. At the condensation scale, level bifurcation is a candidate mechanism. If its three postulates hold, it gives $B=-L$, charge neutrality and the observed composition by symmetry, with an efficiency of up to $3.6\times10^{-3}$ per unit bias. Whether it accounts for $\eta_B$ depends on three derivations: a conserved rung label, a globally signed bias, and the transition dynamics.

\section*{Methods and data availability}
All tests were defined in advance in documents whose SHA-256 hashes and UTC timestamps were recorded before execution. Controls that failed their registered thresholds are reported as failures in the repository, together with corrections and their reasons. Networks were integrated with LSODA, DOP853 or Radau at relative tolerances of $10^{-10}$--$10^{-11}$. Exact results were obtained by linear algebra over the rationals or by singular-value decomposition. Code, definitions, registries and outputs are available in the companion repository \cite{repo}, folder \texttt{p14d\_psion\_number\_audit}.

\begin{thebibliography}{99}
\bibitem{planck} Planck Collaboration, Astron.\ Astrophys.\ \textbf{641}, A6 (2020).
\bibitem{sakharov} A.~D.~Sakharov, JETP Lett.\ \textbf{5}, 24 (1967).
\bibitem{krs} V.~A.~Kuzmin, V.~A.~Rubakov, M.~E.~Shaposhnikov, Phys.\ Lett.\ B \textbf{155}, 36 (1985).
\bibitem{fy} M.~Fukugita, T.~Yanagida, Phys.\ Lett.\ B \textbf{174}, 45 (1986).
\bibitem{bdp} W.~Buchm\"uller, P.~Di~Bari, M.~Pl\"umacher, Ann.\ Phys.\ \textbf{315}, 305 (2005).
\bibitem{ad} I.~Affleck, M.~Dine, Nucl.\ Phys.\ B \textbf{249}, 361 (1985).
\bibitem{phf} K.~M.~Maia, \emph{The Psion Hypothesis: Foundations}, v85 (2026).
\bibitem{repo} K.~M.~Maia, companion repository \texttt{psion-proofs}, \url{https://github.com/Maia-KarlosMarden/psion-proofs}.
\bibitem{witten} E.~Witten, Nucl.\ Phys.\ B \textbf{223}, 433 (1983).
\bibitem{sk} Super-Kamiokande Collaboration, search for nucleon decay via $n\to\bar\nu\pi^0$ and $p\to\bar\nu\pi^+$, arXiv:1305.4391.
\bibitem{kamland} KamLAND Collaboration, Phys.\ Rev.\ Lett.\ \textbf{96}, 101802 (2006), hep-ex/0512059.
\bibitem{cdg} A.~G.~Cohen, A.~De~R\'ujula, S.~L.~Glashow, Astrophys.\ J.\ \textbf{495}, 539 (1998).
\bibitem{kw} E.~W.~Kolb, S.~Wolfram, Nucl.\ Phys.\ B \textbf{172}, 224 (1980).
\bibitem{jr} R.~Jackiw, C.~Rebbi, Phys.\ Rev.\ D \textbf{13}, 3398 (1976).
\end{thebibliography}

\end{document}
